\section{Introduction}

We study how faithfully distances between strings can be represented by
distances in $\ell_1$. The distance on strings is ordinary edit distance:
$\ED(x,y)$ is the least number of single-symbol insertions, deletions and
substitutions that transform $x$ into $y$, with every operation costing
one. An insertion or deletion changes the positions of an entire suffix.
As a result, the distance depends on the best alignment of the two strings,
and a useful representation must account for alignments over all scales.

For a finite alphabet $\Sigma$ and an integer $d\ge1$, let
$\Sigma^{\le d}$ be the set of all strings of length at most $d$, including
the empty string. The cap restricts the endpoints; intermediate strings
in an edit script may have any length. Write $\ell_1$ for the real space of
absolutely summable sequences and define
\[
 E_\Sigma(d)=
 \inf_{f:\Sigma^{\le d}\hookrightarrow\ell_1}
 \left(\max_{x\ne y}\frac{\norm{f(x)-f(y)}}{\ED(x,y)}\right)
 \left(\max_{x\ne y}\frac{\ED(x,y)}{\norm{f(x)-f(y)}}\right).
\]
The infimum is over injective maps. It imposes no restriction on target
dimension or on the cost of computing the map. Equivalently, an embedding
has distortion at most $D$ when it can be rescaled so that every image
distance lies between $\ED(x,y)/D$ and $\ED(x,y)$.

\Needspace{11\baselineskip}
\begin{theorem}\label{thm:main}
There are absolute constants $c,C>0$ and an integer $d_0$ such that, for
every integer $d\ge d_0$ and every finite alphabet $\Sigma$ with
$|\Sigma|\ge2$,
\begin{equation}\label{eq:main}
 \exp\!\bigl(c\sqrt{\log d\,\log\log d}\bigr)
 \le E_\Sigma(d)\le
 \exp\!\bigl(C\sqrt{\log d\,\log\log d}\bigr).
\end{equation}
The lower bound is witnessed by a subset of binary strings of one common
length at most $d$. The same bounds hold for
$\sup_{\Sigma:\,2\le|\Sigma|<\infty}E_\Sigma(d)$.
\end{theorem}

Logarithms are natural unless a base is specified. The constants and
threshold in the theorem are independent of the alphabet, including when
its size grows with $d$. The conclusion determines the order of
$\log E_\Sigma(d)$; the constants in the exponent remain unspecified.

\subsection*{The earlier bounds}

The single-symbol error model is classical. Levenshtein studied codes
correcting insertions, deletions and substitutions
\cite{Levenshtein1966}, and Wagner and Fischer described edit scripts
through increasing traces \cite{WagnerFischer1974}. Their alignment
language will be useful below: any lower bound must control every
increasing matching of equal symbols, including matchings that cross the
boundaries used to construct the strings.

Ostrovsky and Rabani \cite{OstrovskyRabani2007} proved the upper scale in
Theorem~\ref{thm:main} for fixed-length binary strings by recursively
comparing collections of overlapping substrings. They also noted
extensions to larger alphabets and varying lengths. We use their
fixed-length theorem and give complete reductions with constants uniform
over all finite alphabets and with the empty string included.

For lower bounds, Andoni, Deza, Gupta, Indyk and Raskhodnikova
\cite{AndoniDezaGuptaIndykRaskhodnikova2003} constructed binary subsets
whose $L_1$ distortion approaches $3/2$. Khot and Naor
\cite[Theorem~4.1]{KhotNaor2006} obtained an
$\Omega(\sqrt{\log d/\log\log d})$ lower bound by
studying cuts and Fourier analysis of noisy shifts. Their
insertion--deletion metric is within a factor of two of the convention
used here. Krauthgamer and Rabani
\cite[Theorem~1.3]{KrauthgamerRabani2009} then proved an
$\Omega(\log d)$ lower bound for binary strings of length $d$.
The construction below reaches the scale of the Ostrovsky--Rabani
exponent.

\subsection*{How the lower bound works}

The construction joins a geometric fact about strings to an analytic
fact about $\ell_1$. At the geometric level, a word is a long list of
rows. Each row contains a payload made of lower-level words, followed by
a long marker consisting of a fresh symbol. The payload has fixed slots;
each slot has its own alphabet tag, so different slots cannot match.
The lower-level words depend on states. In each component, repeatedly
adding one distinguished step cycles the state with a prime period;
different components use different primes.
Advancing every component by one state step shifts the row list by one,
which costs only the deletion and insertion of one row.

Advancing the distinguished component behaves differently. We repeat
that component in many separately tagged slots. For any proposed pair
of rows, either those repeated slots remain separated or the row shift
that aligns them leaves many of the other prime-period components
separated. A global alignment can still send one source payload into
several target rows. In that event, the matching skips a complete target
marker: the source payload contains no copy of the marker symbol.
The skipped markers for different payloads are disjoint. This turns
the row comparison into a lower bound for arbitrary alignments.

The analytic comparison concerns the average image distance produced
by a group displacement. Proposition~\ref{prop:displacement} bounds
the displacement in the distinguished component by the simultaneous
step and the average component steps. The proof expresses a finite
$\ell_1$ metric as a nonnegative sum of cut metrics, then expands each
cut in Fourier characters. A character involving few components is
detected by the simultaneous step: distinct prime denominators prevent
its phase from being an integer. A character involving many components
is detected by the sum of the component averages.

Relative to word length, the guaranteed pointwise separation decreases
by a fixed factor at each level. Under the same normalization, the
analytic inequality forces the allowed averaged image displacement
to decrease much faster: after $k$ levels, the ratio between these two
bounds is $\exp(\Omega(k\log k))$. Meanwhile, the logarithmic word
length is $O(k^2\log k)$. A delimiter code transfers the words to binary
with a loss polynomial in $k$, and choosing $k$ directly from the cap $d$
supplies a witness for every sufficiently large $d$.

Section~\ref{sec:displacement} proves the alignment facts and the
prime-product displacement inequality. Section~\ref{sec:construction}
builds the words and proves their separation and analytic recurrence
together. Section~\ref{sec:binary} supplies the delimiter at the
binary transfer and completes the parameter calculation.
Section~\ref{sec:upper} derives the uniform upper bound by hashing the
alphabet, padding to one binary length, and averaging over all hash maps.

The companion manuscripts on circle constructions \cite{OpenAICircle2026}
and tree constructions \cite{OpenAITree2026} give independent lower-bound
mechanisms and additional binary coding results. No result from a companion
manuscript is used; the upper bound uses the stated external theorem.
